
\subsection*{2.9. Pointed cocommutative Hopf algebras over a field}
\phantomsection
\label{sga3:VIIB:2.9}
\addcontentsline{toc}{subsection}{\SGARefLink{sga3:VIIB:2.9}{2.9}. Pointed cocommutative Hopf algebras over a field}
\footnote{\label{sga3:VIIB:nde:118}The editors added this subsection
to express Proposition 2.5.2, together with the variant noted in
editorial note 100, in the language of cocommutative Hopf
algebras.}

Let \(k\) be a field, let \(C\) be a \(k\)-coalgebra, and endow
\[
  A=C^*
\]
with the structure of a profinite \(k\)-algebra, not necessarily
commutative, described in \SGARefLink{sga3:VIIB:2.8}{2.8}. By \SGARefLink{sga3:VIIB:0.2.2}{0.2.2},
\[
  C=A^\dagger=\operatorname{Hom}_c(A,k).
\]
The map
\[
  D\longmapsto
  D^\perp=\{f\in A\mid f(D)=0\}
\]
is therefore a bijection from the subcoalgebras of \(C\) to the
closed two-sided ideals of \(A\) \textup{(which we shall simply
call ideals below)}. Its inverse is
\[
  I\longmapsto
  I^\perp
  =
  \{x\in C=A^\dagger\mid x(I)=0\}.
\]
Since every maximal closed ideal is a maximal open ideal
\textup{(cf.\ \SGARefLink{sga3:VIIB:0.2.1}{0.2.1})}, every subcoalgebra contains a simple
subcoalgebra, necessarily finite-dimensional.

Recall that a subcoalgebra \(D\subset C\) is called
\emph{irreducible} if it contains a unique simple subcoalgebra
\(S_0\). Equivalently,
\[
  \mathfrak m_0=S_0^\perp
\]
is the unique open maximal ideal containing \(D^\perp\), or,
equivalently,
\[
  D^\perp+\mathfrak m=A
  \qquad
  \text{for every open maximal ideal }
  \mathfrak m\ne\mathfrak m_0.
\]
This applies in particular when \(D=S_0\).

Let \(\Sigma_0\) be the sum of all irreducible subcoalgebras
\(C_i\) containing \(S_0\). It is plainly a subcoalgebra, and it is
irreducible because, for every \(\mathfrak m\ne\mathfrak m_0\),
\SGARefLink{sga3:VIIB:0.2.D}{0.2.D} gives
\[
  \mathfrak m+\bigcap_i C_i^\perp
  =
  \bigcap_i\bigl(\mathfrak m+C_i^\perp\bigr)
  =
  A.
\]
The coalgebra \(\Sigma_0\) is called the \emph{irreducible
component of \(C\) corresponding to \(S_0\)}.
\textit{The French prints \(C_0\) in this last phrase; the
preceding definition makes \(S_0\) the intended simple
subcoalgebra.}

A coalgebra \(C\) is called \emph{pointed} if every simple
\(k\)-subcoalgebra of \(C\) is one-dimensional. Equivalently,
\[
  A/\mathfrak m=k
\]
for every open maximal ideal \(\mathfrak m\) of \(A\). Recall also
that \(C\) is called \emph{connected} if it is irreducible and
pointed. If \(C\) is a bialgebra, then it is connected if and only
if it is irreducible, since \(k\cdot1_C\) is a simple
subcoalgebra.

Suppose from now on that \(C\) is cocommutative. Then \(A\) is
commutative and hence is the product of its local components
\[
  A_{\mathfrak m},
  \qquad
  \mathfrak m\in\Upsilon(A);
\]
cf.\ \SGARefLink{sga3:VIIB:0.1.1}{0.1.1}. Put
\[
  S_{\mathfrak m}
  =
  \mathfrak m^\perp
  \simeq
  (A/\mathfrak m)^*,
\]
and let \(\Sigma_{\mathfrak m}\) be its irreducible component.
This component can be described explicitly. Let
\[
  J_{\mathfrak m}
  =
  \ker\bigl(A\longrightarrow A_{\mathfrak m}\bigr).
\]
Then \(J_{\mathfrak m}\subset\mathfrak m\), and it is the smallest
closed ideal \(I\subset\mathfrak m\) satisfying
\[
  I+\mathfrak m'=A
  \qquad
  \text{for every open maximal ideal }
  \mathfrak m'\ne\mathfrak m.
\]
Indeed, any such \(I\) contains \(A_{\mathfrak m'}\) for every
\(\mathfrak m'\ne\mathfrak m\), and hence contains
\(J_{\mathfrak m}\). Since
\[
  A=J_{\mathfrak m}\oplus A_{\mathfrak m},
\]
we obtain
\[
  \Sigma_{\mathfrak m}
  =
  J_{\mathfrak m}^{\perp}
  \simeq
  A_{\mathfrak m}^{\dagger}.
\]
We can now prove:

\medskip
\noindent\textbf{Proposition 2.9.1.}
\phantomsection
\label{sga3:VIIB:2.9.1}
{\itshape Let \(k\) be a field.
\begin{enumerate}
\item[\textup{\SGAInlineRefTarget{sga3:VIIB:2.9.1:item:i}(i)}]
Let \(G\) be a formal \(k\)-group all of whose residue fields are
equal to \(k\). Then \(G_e\) is the constant \(k\)-group \(M_k\),
where
\[
  M
  =
  G(k)
  =
  \left\{
    x\in\mathbf H(G)
    \,\middle|\,
    \varepsilon(x)=1,\quad
    \Delta(x)=x\otimes x
  \right\},
\]
and \(\mathbf H(G)\) is the semidirect product of
\(\mathbf H(G_0)\) by the group algebra \(kM\)
\textup{(cf.\ \SGARefLink{sga3:VIIB:2.4.B}{2.4.B})}.

\item[\textup{(ii)}]
Equivalently, let \(H\) be a pointed cocommutative Hopf
\(k\)-algebra. Then \(H\) is the semidirect product of the
irreducible component \(H_0\) of the unit \(1_H\) by \(kM\), where
\[
  M
  =
  \{x\in H\mid \varepsilon(x)=1,\ \Delta(x)=x\otimes x\}.
\]
\end{enumerate}}

\textit{Proof.}
We first prove \textup{\SGARefLink{sga3:VIIB:2.9.1:item:i}{(i)}}. Since every residue field of \(G\) is
equal to \(k\), the projection
\[
  \pi:G\longrightarrow G_e
\]
has a section \(s:G_e\to G\), defined at every point \(g\in G\) by
\[
  \mathcal O_{G,g}\longrightarrow\kappa(g)=k.
\]
Moreover, for all \(g,h\in G\), the algebra
\[
  \mathcal O_{G,g}
  \widehat\otimes_k
  \mathcal O_{G,h}
\]
is local with residue field \(k\). Thus \(s\times s\) is a section
of
\[
  \pi\times\pi:
  G\times G\longrightarrow
  (G\times G)_e=G_e\times G_e.
\]
Since \(\pi\) is a group morphism,
\[
  \pi\circ m_G\circ(s\times s)
  =
  m_{G_e}
  =
  \pi\circ s\circ m_{G_e}.
\]
As \(\pi\) is an epimorphism, it follows that \(s\) is a group
morphism. Hence
\[
  G=G_0\rtimes G_e,
\]
and \(\mathbf H(G)\) is the semidirect product of
\(\mathbf H(G_0)\) by \(\mathbf H(G_e)\).

Since every residue field of \(G\) is \(k\),
\(\mathbf H(G_e)\) is the group algebra \(kM\), where
\[
  M=G_e(k);
\]
cf.\ \SGARefLink{sga3:VIIB:2.5.A}{2.5.A}. Finally, because \(G_0\) is infinitesimal, the map
\[
  G(k)\longrightarrow G_e(k)
\]
is injective. It is also surjective, since it has the section
induced by \(s\). Thus \(M=G(k)\), proving \textup{\SGARefLink{sga3:vii:paragraph:5-1:equation:i}{(i)}}.

For \textup{\SGARefLink{sga3:vii:paragraph:5-1:equation:ii}{(ii)}}, it remains only to show that
\[
  H_0=\mathbf H(G_0).
\]
The unit \(1_H\) of \(H\) is the augmentation
\[
  \varepsilon_{\mathcal A}:\mathcal A\longrightarrow k,
\]
which is the identity section of \(G(k)\). Therefore the local
component of the affine algebra \(\mathcal A(G)\) corresponding to
\(H_0\) is precisely \(\mathcal A(G_0)\). By the description above,
\[
  H_0
  =
  \mathcal A(G_0)^\dagger
  =
  \mathbf H(G_0).
\]
This proves the proposition.

\medskip
\noindent\textbf{Remarks 2.9.2.}
\phantomsection
\label{sga3:VIIB:2.9.2}
\begin{enumerate}
\item[\textup{(a)}]
The preceding proposition, which is implicit in \SGARefLink{sga3:VIIB:2.5.2:proposition}{2.5.2}, was obtained
independently by B.\ Kostant \textup{(cf.\ [Sw69], Preface)}.
Combined with Cartier's \SGARefLink{sga3:VIIB:3.3:cartier-theorem}{Theorem 3.3} below
\textup{(cf.\ [Ca62], \S~12, Theorem 3)}, which was also obtained
by Kostant \textup{(cf.\ [Sw69], loc.\ cit.)}, this result is often
called the \emph{Cartier--Gabriel--Kostant theorem}.

\item[\textup{(b)}]
In the form \textup{\SGARefLink{sga3:vii:paragraph:5-1:equation:ii}{(ii)}}, \SGARefLink{sga3:VIIB:2.9.1}{Proposition 2.9.1} was extended by
R.\ G.\ Heyneman and M.\ E.\ Sweedler to the case in which \(H\)
is assumed to be pointed and the direct sum of its irreducible
components, but is not necessarily cocommutative; see [HS69],
Theorem 3.5.8.
\end{enumerate}
